\subsection{Exact characteristic-two plane monodromy in every degree}
\label{subsec:char2-uniform}

\begin{lemma}[Two elementary group-theoretic inputs]
\label{lem:char2-group-inputs}
For every integer $a\ge8$, there is an odd prime $\ell_a$ with
$\operatorname{ord}_{\ell_a}(2)=a$.
For $N\ge4$, the first cohomology of $\mathrm{GL}_N(\mathbf F_2)$
with coefficients in its natural module or its dual is zero.
\end{lemma}
\begin{proof}
We include a proof of the needed special case of the primitive-divisor
theorem. If an odd prime $\ell$ divides $\Phi_a(2)$ without having
order $a$, the cyclotomic factorization and the lifting-the-exponent
identity show that
$a=\operatorname{ord}_{\ell}(2)\ell^j$ with $j\ge1$, and that
$v_\ell(\Phi_a(2))=1$. Indeed, substitute
$v_\ell(2^{fm}-1)=v_\ell(2^f-1)+v_\ell(m)$, where
$f=\operatorname{ord}_\ell(2)$, in the M\"obius formula for
$v_\ell(\Phi_a(2))$; the sum is zero unless $a/f$ is a power of
$\ell$, when it is one. Since $f\mid\ell-1$, such an $\ell$ is the
largest prime factor $P$ of $a$. Thus, if no primitive prime occurs,
$\Phi_a(2)$ is either $1$ or $P$.
On the other hand the cyclotomic product gives
\[
 \Phi_a(2)>2^{\varphi(a)}\prod_{j\ge1}(1-2^{-j})
 >2^{\varphi(a)-2}.
\]
For the last bound, the first three factors have product $21/64$,
and the remaining product is at least
$1-\sum_{j\ge4}2^{-j}=7/8$; their product exceeds $1/4$.
If $P\ge7$, then $\varphi(a)\ge P-1$ and $2^{P-3}>P$.
If $P\le3$ and $a\ge8$, then $\varphi(a)\ge4$ and the same
conclusion holds. If $P=5$, either $\varphi(a)\ge8$ or $a=10$;
the exceptional calculation is $\Phi_{10}(2)=11$.
This contradiction proves the first assertion. The cyclotomic
valuation calculation also agrees with
\cite[Lemma~3.1(b)--(c), pp.~125--126]{GlasbyEtAl2017}, whose
statement and proof were opened.

For the second assertion let $z$ be a $1$-cocycle with values in
$\mathbf F_2^N$ and put $t_{ij}=1+E_{ij}$.
The relation $t_{ij}^2=1$ forces the $j$th coordinate of
$z(t_{ij})$ to vanish. For $l\ne i,j$, choose
$k\notin\{i,j,l\}$. Comparing the cocycle on the commuting pair
$t_{ij},t_{kl}$ shows that its $l$th coordinate also vanishes.
Thus $z(t_{ij})=a_{ij}e_i$.
For distinct $i,j,k$, apply the cocycle to
$[t_{ij},t_{jk}]=t_{ik}$; it gives $a_{ik}=a_{jk}$.
Consequently $a_{ij}=b_j$ for a single vector $b$, and
$z(t_{ij})=(t_{ij}-1)b$. Since these transvections generate the
group, $z$ is a coboundary. Composing with inverse transpose
gives the dual-module assertion.
\end{proof}

\begin{theorem}[All plane degrees in characteristic two]
\label{thm:plane-char2-all-groups}
Let $k$ be any perfect field of characteristic two, let $U_d$ be
the ordinary smooth locus in the projective space of degree-$d$
plane forms, and choose a geometric generic base point $\bar\eta$.
Write
\[
 g=(d-1)(d-2)/2,\qquad
 V=H^1_{\mathrm{et}}(C_{\bar\eta},\mathbf F_2).
\]
For $d\ge4$, the images of $\pi_1(U_d\otimes_k\bar k,\bar\eta)$
and $\pi_1(U_d,\bar\eta)$ on $V$ are equal and are
\[
 \begin{cases}
 \mathrm{GL}_g(\mathbf F_2),&d\text{ even},\\
 \operatorname{Stab}(\ell)
   \simeq(\mathbf F_2)^{g-1}\rtimes
                  \mathrm{GL}_{g-1}(\mathbf F_2),&d\text{ odd},
 \end{cases}
\]
where $\ell\ne0$ is the canonical covector determined by
$\mathcal O_C((d-3)/2)\otimes B_C^{-1}$ in odd degree.
The nonzero representation is irreducible exactly in even degree;
its determinant is always trivial. In degrees one and two $V=0$;
in degree three it has rank one and trivial, hence irreducible,
monodromy. Thus there is no threshold after which all plane
degrees have full linear or irreducible monodromy.
\end{theorem}
\begin{proof}
The degrees at most four, the quintic stabilizer, and the sextic
full group have been proved above. We induct on $d\ge7$.
Work first over $\bar k$ and on the dual space $W_0=V^\vee$.
Theorem~\ref{thm:all-char2-theta-orbits} proves transitivity on
$W_0\setminus\{0\}$ when $d$ is even and on
$W_0\setminus\{0,\ell\}$ when $d$ is odd.
In the odd case put $W=W_0/\langle\ell\rangle$; in the even
case put $W=W_0$. The induced group on $W$ is therefore transitive
on its nonzero vectors. Its dimension is
\[
 N=\begin{cases}g,&d\text{ even},\\g-1,&d\text{ odd}.\end{cases}
\]

We specify the boundary subgroup used in the induction. The
universal plane-curve coherent $H^1$ is a vector bundle of rank
$g$, including singular fibres. The locus where its Frobenius is
bijective is open, and the Lang construction gives an etale
$\mathbf F_2$-local system on this whole locus. It contains the
nodal curves $C'\cup L$, where $C'$ is ordinary smooth of degree
$d-1$ and $L$ meets it transversely in $d-1$ distinct points.
The smooth ordinary locus is dense in this normal parameter
space. Restriction of connected finite etale covers to a dense
open remains connected, so its fundamental group maps onto the
fundamental group of the larger locus. Boundary monodromy thus
lies in the group being determined.

The normalization sequence for $C'\cup L$ gives
\[
 0\longrightarrow H^1(\Gamma,\mathbf F_2)
 \longrightarrow H^1_{\mathrm{et}}(C'\cup L,\mathbf F_2)
 \longrightarrow H^1_{\mathrm{et}}(C',\mathbf F_2)
 \longrightarrow0.
\]
Here the graph space has dimension $d-2$ and its monodromy factors
through the permutation group $S_{d-1}$ of the nodes. The
projection from the parameter space of transverse pairs $(C',L)$
to that of $C'$ is smooth, surjective, and has geometrically
connected fibres, namely nonempty opens of the dual plane.
It therefore surjects on fundamental groups and supplies the
entire already determined monodromy of $C'$.

Write $r=(d-2)(d-3)/2$. If $d$ is even, the preceding odd-degree
group has a quotient $\mathrm{GL}_{r-1}(2)$; the kernel of this
quotient in the boundary group has prime divisors only $2$ and
those dividing $(d-1)!$. If $d$ is odd, the dual boundary
sequence has an invariant subspace of dimension $r$ with full
$\mathrm{GL}_r(2)$ action. The nonzero invariant vector $\ell$
cannot lie in that subspace, since it has no nonzero invariant
vector. It remains nonzero at the boundary: an invariant section
of a finite etale sheaf on the dense open extends uniquely over
the normal larger locus, and the zero and nonzero sections are
disjoint. Thus the rank-$r$ subspace injects into $W$, with quotient
a subquotient of the graph permutation space. In this case also,
the kernel of the boundary quotient onto $\mathrm{GL}_r(2)$ has
prime divisors only $2$ and those dividing $(d-1)!$. These kernel
claims follow by writing the representation in block triangular
form: its off-diagonal part is a $2$-group, the preceding odd-degree
unipotent radical is a $2$-group, and its remaining diagonal
action factors through $S_{d-1}$.

Set $a=\lfloor N/2\rfloor+1$. The elementary inequalities
\[
 8\le a<a+1<N-1,\qquad a>d-1,\qquad
 a+1\le\begin{cases}r-1,&d\text{ even},\\r,&d\text{ odd}
 \end{cases}
\]
hold for even $d\ge8$ and odd $d\ge7$. For example the two initial
cases have $(N,r,a)=(21,15,11)$ and $(14,10,8)$.
For the upper bound, using $\lfloor N/2\rfloor\le N/2$
reduces it to $g\ge2d+2$ in even degree and $g\ge2d-1$ in odd
degree; both hold at the indicated initial degree and their
left-minus-right differences increase. The other inequalities
follow directly from the same genus formula.
For $j=a,a+1$, choose the prime $\ell_j$ in
Lemma~\ref{lem:char2-group-inputs}. We have $\ell_j>j>d-1$.
Also $\ell_j\nmid N$: otherwise $\ell_j>N/2$ forces $N=\ell_j$,
whereas $j\mid N-1$ and $N/2<j<N-1$, an impossibility.

In the boundary general linear quotient choose an element of
order $\ell_j$ with one irreducible $j$-dimensional block and
identity on its complement. The kernel has order prime to
$\ell_j$, so a Sylow subgroup of its inverse image supplies an
element of order $\ell_j$ with that image. It acts trivially on
the graph quotient, because $\ell_j\nmid(d-1)!$.
Semisimplicity of this odd-order action shows that its
characteristic polynomial on $W$ is
\[
                    (T-1)^{N-j}f_j(T),
 \qquad f_j\text{ irreducible of degree }j.
\]

Apply the opened classification in
\cite[Lemma~3.3, p.~6]{ChurikovVasilevZvezdina2022}.
Every proper extension-field alternative lies inside
$\Gamma\mathrm L_{N/b}(2^b)$ for some divisor $b>1$ of $N$.
Our order-$\ell_j$ elements belong to its linear subgroup,
since $\ell_j\nmid b$. Their fixed-space dimensions are therefore
divisible by $b$. The consecutive dimensions $N-a$ and $N-a-1$
force $b=1$, excluding all these alternatives, including the
one-dimensional semilinear case. Over the original field the
remaining symplectic alternative lies in $\mathrm{Sp}_N(2)$:
its normalizer preserves the unique invariant alternating form,
and the only nonzero scalar is one. One of $a,a+1$ is odd.
For this $j$, the irreducible factor $f_j$ is not reciprocal,
since inversion is not in the odd-order subgroup generated by
$2$ modulo $\ell_j$. Its characteristic polynomial contradicts
preservation of a nondegenerate alternating form. The $G_2$
alternative over a larger field was already excluded, and its
field-two exceptions have dimension six. The other exceptional
field-two cases have dimension four. Here $N\ge14$.
The classification therefore forces the image on $W$ to be
$\mathrm{SL}_N(2)=\mathrm{GL}_N(2)$.

For even $d$ this finishes the geometric assertion. For odd $d$
the kernel of the image in $\mathrm{GL}(W)$ is an invariant
subgroup of $\operatorname{Hom}(W,\langle\ell\rangle)$.
This is the irreducible dual natural module, so the kernel is
zero or the whole module. If it were zero, the group would be a
complement. Lemma~\ref{lem:char2-group-inputs} makes every such
complement conjugate to the block diagonal group. That group has
two separate orbits on $W_0\setminus\{0,\ell\}$, contradicting
the theta-orbit theorem. Thus the kernel is the whole module,
giving the full stabilizer and completing the induction.

The distinguished odd-degree covector is defined over the original
perfect field, as proved in the explicit Hasse--Witt construction.
Arithmetic monodromy is contained in its stabilizer and contains
the geometric group just determined; in even degree it is
contained in the full general linear group. Thus the two images
are equal in both cases. Irreducibility and the determinant
assertions follow immediately; $\mathbf F_2^\times$ is trivial.
\end{proof}
